\documentclass[12pt,a4paper]{ctexart}
\usepackage{amsmath,amssymb}
\usepackage{booktabs}
\usepackage{array}
\usepackage{geometry}
\usepackage{enumitem}
\usepackage[hidelinks]{hyperref}
\geometry{margin=2.5cm}
\linespread{1.25}
\setlength{\parskip}{0.4em}
\setlength{\parindent}{2em}
\renewcommand{\arraystretch}{1.15}
\ctexset{section={format={\heiti \zihao{3}}}, subsection={format={\heiti \zihao{4}}}}
\setcounter{secnumdepth}{-1}

\begin{document}

\begin{center}
{\heiti\zihao{3} 2026 CUMCM A 题论文模拟评阅意见\par}
\vspace{0.4em}
{\zihao{-4} （评委视角 · 模拟评阅 · 基于 2026-09-13 定稿版 · 25 页）\par}
\end{center}

\section{一、总评与得分}

\textbf{总评：82/100，较上一轮外部审稿（76/100）有实质提升。} 已消除的硬伤包括：$C_a$ 物理口径统一（5.1 节声明为药材等效平衡水分浓度）、移动边界守恒推导补齐（5.4 节从干固体与水分守恒律出发）、「严格守恒」与「离散通量严格闭合」的措辞区分、表面传质主导的错误判断（6.3 节改为 Bi$_m$ 数据支撑的「共同控制」）、单因素可叠加性经最不利联合工况实算验证（7.7 节）。当前版本处于\textbf{「国二较稳、具备国一讨论价值」}区间；与国一之间最主要的差距不在建模与计算，而在\textbf{图形化呈现与摘要可读性}。

\begin{center}
{\zihao{6}\begin{tabular}{lccp{5.8cm}}
\toprule
维度 & 得分 & 满分 & 评语要点 \\
\midrule
问题理解 & 9 & 10 & 四问条件差异梳理清楚，2.4 节对照表降低错用参数风险 \\
模型合理性 & 17 & 20 & 有效参数观点（H5）与潜热处理逻辑自洽；数据自洽性可再补一句 \\
数值求解 & 19 & 20 & 统一守恒框架、真实表面节点、动边界随体坐标，格式与误差分析专业 \\
结果与检验 & 14 & 15 & 五重验证+消融+灵敏度+联合工况；物理水量核验与二维对照仍缺 \\
创新辨识度 & 11 & 15 & 2$\times$2 消融与随体坐标推导是亮点，但表达未突出到评委一眼可见 \\
写作与呈现 & 8 & 15 & 无图是最大短板；摘要过密；数值细节占比偏高 \\
规范与复现 & 4 & 5 & 基本合规；源代码附录与 AI 声明句待补 \\
\midrule
合计 & 82 & 100 & 国二较稳，补图后具备国一讨论价值 \\
\bottomrule
\end{tabular}}
\end{center}

\section{二、主要优势（评语与证据）}

\begin{enumerate}[label=(\arabic*),leftmargin=2em]
\item \textbf{数值框架统一且严谨。}表面节点置于真实表面、Robin 通量与内部界面通量逐层对消，离散闭合差达 $10^{-7}$（问题 1 机器精度 $3.9\times10^{-19}$，7.3 节）；Crank--Nicolson 与交错 Picard 的收敛判据、停止阈值的精度依据（5.5 节）都有交代。这类「格式层面想清楚再动手」的素养在竞赛论文中属上乘。
\item \textbf{移动边界处理有推导深度。}5.4 节从干固体质量守恒 $\mathrm{d}M_d=\rho_d\mathrm{d}V$ 与水分守恒出发，利用 $\partial\rho_d/\partial t+\nabla\cdot(\rho_d v)=0$ 严格消去对流项，再以假设 H5 说明 $\rho_d$ 均匀近似的归属。多数参赛队直接套用随体坐标公式，本文的推导层次明显更高，且规避了审稿报告中指出的逻辑缺口。
\item \textbf{验证体系完整且互证。}Bessel 解析解四位小数一致（7.1）、空间/时间收敛链（7.2）、离散守恒（7.3）、长时渐近（7.4）、活化能文献互证（7.5）五重检验；问题 3、4 针对插值取值导致的降阶不硬套理想二阶、改用实测收敛阶估计残余误差（$t^*$ 链外推 57.1732/50.8239 h，不确定度 $\pm0.03$ h/$\pm0.01$ h），误差意识显著优于同类论文。
\item \textbf{表 7 的物性$\times$几何消融是全文创新中心。}用同一套验证离散把「物性变化」与「收缩变化」分离为 $+126\%$ 与 $-56\%\sim-61\%$ 两个机制分量，再合成净效应 $-11\%$，把简单结果对比升级为机制分析；6.3 节用 Bi$_m$（2.8$\to$2.4$\to$15.5$\to$超 200）这一有量纲数支撑「共同控制」的判断，说服力强。
\item \textbf{敏感性与稳健性分析闭环。}8 个参数变体+最不利联合工况实算（68.22 h，仍在「2--3 天」范围），直接回应了非线性模型下「单因素不可加」的质疑，结论稳健性有实算支撑而非声明。
\item \textbf{潜热的处理克制而有深度。}7.6 节给出 $q_{\mathrm{evap}}/q_{\mathrm{sens}}\approx3\sim5$ 的数量级估计与 $+40\%\sim50\%$ 的时长敏感性区间，既说明显式叠加潜热项的风险（与经验参数重复计账），又不越界声称「已被吸收」——分寸感符合专业写作规范。
\item \textbf{结论直接回应四问。}8.3 节集中给出四问数值答案、消融结论与工艺启示（温度控制是核心），评委快评时可在此页拿全关键信息。
\end{enumerate}

\section{三、主要问题与修改建议（按优先级）}

\subsection{P0（合规风险，建议优先处理）}

\textbf{问题 1：附录未包含完整源程序代码。}2026 格式规范第五条明确要求「论文附录内容应包括支撑材料的文件列表，建模所用到的全部完整、可运行的源程序代码」，并警示「缺少必要的源程序……可能会被取消评奖资格」。当前附录仅有文件列表（八个程序合计约 1345 行，均未印入）。支撑材料压缩包中的代码不能替代附录位置的要求。\textbf{建议：}将八个源程序以等宽小字号排版补入附录（附录页数不限，预计增加约 20 页）；若赛区有不同口径，以赛区书面要求为准。

\textbf{问题 2：AI 使用声明缺少规定格式的标准句。}2026 年 AI 工具使用规定要求在论文参考文献之前设置声明，形式为「本参赛队在竞赛过程中使用了AI工具，主要用于【简要用途，如语言润色、代码调试等】，详细使用情况见支撑材料」。当前论文为详情附录（内容四项齐全、位置正确），但缺少这句规定格式的简要声明。同时支撑材料中应包含固定文件名「AI工具使用详情.pdf」（论文附录的详情内容导出为独立 PDF 即可，四项内容已覆盖）。\textbf{建议：}在附录标题下补一句标准声明；文件列表补「AI工具使用详情.pdf」一项。

\subsection{P1（呈现短板，直接决定快评印象）}

\textbf{问题 3：全文无图，是当前最大的呈现短板。}25 页正文 8 表 0 图。评委在快速评阅中只能通过表格与文字寻找主结论，辨识度受损。至少应补 4 幅：
\begin{enumerate}[label=(\alph*),leftmargin=2em]
\item \textbf{技术路线图}：数据输入$\to$固定/变物性模型$\to$固定/移动边界$\to$数值求解$\to$五重验证$\to$结论；
\item \textbf{问题 1 径向剖面}：表 1、表 2 数据直接可画，2--3 个时刻的温度与水分浓度径向剖面（水分边界层约 5 mm 的结论一图即明）；
\item \textbf{问题 3/4 含水率时程}：中心与表面含水率曲线+$C=0.15$ 阈值线与 $t^*$ 标注（问题 3、4 对比）；
\item \textbf{表 8 龙卷风图}：八个变体对 $t^*$ 的影响一目了然，且自然突出「$T_a$ 最敏感、$h_m$ 次之」的工艺结论。
\end{enumerate}
图可直接用现有求解器输出数据绘制，不引入任何新计算。

\textbf{问题 4：摘要过密。}一页内挤入四问、数值方法、四重检验与潜热敏感性区间，且摘要中出现三个公式。建议压缩约 1/3：删去摘要中的 $D(C,T)$ 显式公式与「$+40\%\sim50\%$」细节，四问各按「模型—结果」一句组织，末句概括创新（统一守恒框架+移动边界+消融+多重验证）。

\textbf{问题 5：数值细节占比偏高。}5.6 节逐问加密过程四段约 40 行、7.2 节收敛链数字密集，快评时挤占主线。建议主文只保留各问最终离散与误差结论，完整加密序列与逐格差值移至附录表格。

\subsection{P2（技术增强，冲击国一的加分项）}

\textbf{问题 6：二维对照缺位。}端面效应目前仅由轴向/径向扩散特征时间之比 $(L/(2R))^2\approx39$ 的尺度估计支撑（2.1 节）。建议用较粗离散做一次二维轴对称对照，给出几何中心含水率与烘干时长的一维/二维差值，把「可以忽略」从估计升级为定量结论。

\textbf{问题 7：物理水量核验缺位。}7.3 节诚实地区分了离散闭合与物理守恒（此处理得当），但未给出物理总水量 $\int\rho_d C\,\mathrm{d}V$ 的数值核验。若能按 5.4 节的 $\rho_d$ 定义与收缩 Jacobian 计算全程物理水量变化并对照表面累计通量，可再补一重验证。

\textbf{问题 8：非线性与动边界离散缺制造解。}Bessel 对照仅覆盖常系数线性热传导；非线性水分方程与动边界离散的验证目前依赖收敛链与守恒。建议对 $D(C,T)$ 方程与 $\zeta$ 坐标格式各构造一个制造解（解析余项法），成本低、说服力强。

\textbf{问题 9：题给数据的自洽性说明。}若将附录 4 的 $\rho(C)=760+90C$ 与附件 2 的 $R(t)$ 联合解读，干物质质量在收缩全程存在约 13\% 量级的漂移，即两组题给数据在干物质守恒意义上并非严格自洽。由于 $\rho$ 仅进入热方程（且表 8 显示 $h\pm50\%$ 的影响不足 0.1\%，热场对 $t^*$ 的影响有限），该漂移对水分场与结论无实质影响，但建议在 5.4 节或 7.5 节加一句数据自洽性说明，避免细究型评委的追问。

\section{四、逐节细评}

\begin{center}
{\zihao{5}\begin{tabular}{lp{11.5cm}}
\toprule
章节 & 评语 \\
\midrule
摘要 & B+。四问结论、方法、检验俱全，匿名合规；密度过高，见问题 4 \\
一、问题重述 & A-。四问差异与输出要求完整，2.4 节条件对照表（附录/几何/时段）是本节的加分项 \\
二、模型假设 & A-。H1--H8 逐一给出依据，H5 有效参数观点（含汽化潜热）与 H8 收缩假设表述克制 \\
五、模型建立 & A-。四问统一 PDE 框架，$C_a$ 口径在 5.1 节声明并指向 7.6 节情景分析，动边界推导完整（见优势 2） \\
5.5--5.6 数值方法 & A。格式、截断误差、Picard 判据、步长选取原则专业；篇幅可压缩（问题 5） \\
六、求解结果 & B+。表 1--8 完整、四位小数合规；6.3/6.4 的 Bi$_m$ 机制分析与消融解读是亮点；无图（问题 3） \\
七、模型检验 & A。五重验证+灵敏度+联合工况闭环，7.3 对离散/物理守恒的区分体现学术诚实 \\
八、评价与推广 & A-。8.3 集中回答四问并给出工艺启示；推广措辞克制（只谈可用于品质风险评估、不声称预测成分含量），符合规范 \\
附录 & B。支撑材料文件列表齐全（九个文件+赛题数据副本）；AI 详情内容四项齐全但缺标准声明句；源程序代码待补印 \\
参考文献 & A。17 条全部可追溯（已删除两条无法核验卷页的文献、替换一条），正文引用标注完整 \\
\bottomrule
\end{tabular}}
\end{center}

\section{五、规范合规检查结果（对照 2026 格式规范与 AI 使用规定）}

\begin{center}
{\zihao{6}\begin{tabular}{lp{4.6cm}p{6.2cm}}
\toprule
检查项 & 要求 & 现状 \\
\midrule
摘要页 & 单独一页、含标题关键词、$\le$1 页 & 合规（第 1 页） \\
页码 & 自摘要页起页脚中部连续编号 & 合规 \\
正文页数 & $\le$30 页、无目录 & 合规（正文 20 页） \\
匿名性 & 全文无身份/学校/赛区信息 & 合规（含代码注释与 PDF 元数据） \\
电子版 & PDF 单文件 $\le$20MB、首页为摘要页 & 合规（488KB） \\
支撑材料 & ZIP $\le$20MB、文件列表入附录、与论文一致 & 基本合规（文件列表齐全；打包时注意排除调试脚本与缓存文件） \\
源程序代码 & 附录须包含全部完整可运行代码 & \textbf{待补}（问题 1） \\
AI 声明 & 参考文献前标准声明句+支撑材料详情 PDF & \textbf{待补}（问题 2） \\
\bottomrule
\end{tabular}}
\end{center}

\section{六、若时间有限的执行顺序}

先做 P0（源代码附录、AI 声明句与详情 PDF——涉及资格风险）；再补图 4 幅（投入小、快评收益最大）；随后摘要瘦身与数值细节压缩；P2 三项按剩余时间选做，其中制造解（问题 8）成本最低、二维对照（问题 6）收益最直观。

\vspace{1em}
\noindent\textbf{结语：}本文的计算与验证工作量真实、严谨、可复现，核心结论有实算支撑；上一轮指出的科学性问题已全部消除。当前只需把「最值得评委看到的内容」图形化、把合规两处补齐，即具备送国评的完整面貌。

\end{document}
